%قبل از شروع مسئله ترجیحاً چیزی ننویسید. اما اگر خواستید بنویسید، در قالب لیستی (فقرات و شمارش) بنویسید.
%برای اضافه کردن نمره، بعدی از شروع‌مسئله، [x] را اضافه کنید و به جای x  نمره مورد نظر را قرار دهید. برای زیرمسئله‌ها هم به همین صورت است.
\شروع‌مسئله[20]
%اگر مسئله چند قسمت داشت به این صورت عمل کنید:
%\شروع{زیرمسائل}
%\زیرمسئله[0] % مسئلۀ ۱آ
%\زیرمسئله % مسئلۀ ۱ب
%\زیرمسئله % مسئلۀ ۱پ
%\زیرمسئله % مسئلۀ ۱ت
%\پایان{زیرمسائل}
\متن‌سیاه{مؤلفه‌های قویاً همبند}

در این سؤال باید تعداد و اندازۀ مؤلفه‌های قویاً همبند (یا همبند قوی) یک گراف جهت‌دار را بیابید. برای این کار می‌توانید از الگوریتم \lr{Kosaraju} استفاده کنید. مراحل این الگوریتم به شکل خلاصه عبارتند از 
\begin{enumerate}[1. ]
	\item
ابتدا رأس‌ها را همانند روش مرتب‌سازی توپولوژیک بر حسب عکس زمان پایان \متن‌لاتین{DFS}شان مرتب می‌کنیم (توجه کنید گراف لزوماً \lr{DAG} نیست ولی با همان روش یک ترتیب برای راس‌ها به‌دست می‌آید).
	\item
برحسب ترتیب به دست آمده در مرحلۀ قبل شروع به اجرای
\lr{DFS}
می‌کنیم اما در این مرحله یال‌های گراف را در جهت عکس درنظر می‌گیریم. رأس‌هایی که در این \lr{DFS}
% زده می‌شود تا مولفه‌ای که راس آخر در آن قرار دارد، بدست بیاید. راس‌هایی که 
از آن‌ها عبور می‌کنیم همگی در یک مؤلفۀ قویاً همبند هستند. سپس رأس‌های مؤلفۀ پیدا شده را کنار گذاشته و مرحلۀ $2$ را تکرار می‌کنیم تا همۀ مؤلفه‌های قویاً همبند یکی یکی پیدا شوند.
\end{enumerate}

برای مطالعۀ بیشتر می‌توانید به این
\href{https://www.geeksforgeeks.org/strongly-connected-components}{لینک}
مراجعه کنید.

\متن‌سیاه{ورودی:}\\
در خط اول ورودی $n$ تعداد رأس‌ها ($n \leq 10^6$) و $m$ تعداد یال‌ها ($m \leq 10^6$) آمده است.
در $m$ خط بعدی در هر خط دو عدد $u_i\, v_i$ می‌آیند که نشان دهندۀ وجود یال از رأس $u_i$ به رأس $v_i$ است.
شمارۀ راس‌ها از $0$ تا $n-1$ است.

\متن‌سیاه{خروجی:}\\
در خط اول تعداد مؤلفه‌های قویاً همبندی گراف را چاپ کنید.
در خط دوم اندازۀ مؤلفه های قویاً همبند را به ترتیب اندازه‌شان از کوچک به بزرگ چاپ کنید.

\solution{
پس از اجرای الگوریتم \lr{Kosaraju} به مؤلفه‌های قویاً همبند دست پیدا می‌کنید که با شماردن اندازهٔ مؤلفه‌ها جواب مورد نظر بدست می‌آید.

کد ضمیمه شده برای این سؤال از آقای رضا یاربخش است.

}






%%%%%%%%%%%%%%%%%%%%%%%%%%%%%%%%%%%%%%%%%%%%%%%%%%%%%
% برخی دستورات مفید برای تمرین‌ها. %
%%%%%%%%%%%%%%%%%%%%%%%%%%%%%%%%%%%%%%%%%%%%%%%%%%%%%

% فرمول‌های نمونه:
%$f(x) = \Theta(x)$
%$T(x, y) \leq \log(x) + 2^y + \binom{2n}{n}$
% {\tt code\_function}


% لیست موارد بدون شماره‌گذاری:
%\شروع{فقرات}
%    \فقره مورد اول 
%        \شروع{فقرات}
%            \فقره مورد اول داخلی 
%                \شروع{فقرات}
%                    \فقره اولین مورد خیلی داخلی 
%                    \فقره دومین مورد خیلی داخلی 
%                \پایان{فقرات}
%            \فقره مورد دوم داخلی 
%        \پایان{فقرات}
%    \فقره مورد دوم 
%\پایان{فقرات}

% لیست موارد شماره‌گذاری شده:
%\شروع{شمارش}
%    \فقره یک مورد 
%    \فقره یک مورد دیگر 
%    \فقره باز هم یک مورد دیگر
%\پایان{شمارش}

% برای نوشتن یک فرمول چند خطی که یک شماره بخوره. برای حالت‌های دیگر به سایت زیر مراجعه کنید:
% https://www.overleaf.com/learn/latex/Aligning_equations_with_amsmath
%\begin{align} 
%    \begin{split}
%        (x+y)^3 &= (x+y)^2(x+y) \\
%                &= (x^2+2xy+y^2)(x+y) \\
%                &= (x^3+2x^2y+xy^2) + (x^2y+2xy^2+y^3) \\
%                &= x^3+3x^2y+3xy^2+y^3
%    \end{split}                                 
%\end{align}

% نوشتن ماتریس و آرایه دوبعدی
%\begin{align}
%    A = 
%    \begin{bmatrix}
%        A_{11} & A_{21} \\
%        A_{21} & A_{22}
%    \end{bmatrix}
%\end{align}

% اضافه کردن تصویر با شرح. برای اضافه کردن تصویر بدون شرح، فقط دستور درج‌تصویر را استفاده کنید.
%\شروع{شکل}[h!]
%\تنظیم‌ازوسط
%  \درج‌تصویر[پهنا=0.5\پهنای‌متن]{img.jpg}
%\شرح{شرح تصویر}
%\برچسب{تصویر:۱}
%\پایان{شکل}
%\پایان{زیرمسائل}

%کد زیر را ببینید.
%%\lstinputlisting{source_filename.py}
%\begin{latin}
%\begin{lstlisting}
%# Import the modules
%import sys
%import random
%
%ans = True
%
%while ans:
%question = raw_input("Ask the magic 8 ball a question: (press enter to quit) ")
%
%answers = random.randint(1,8)
%
%if question == "":
%sys.exit()
%
%elif answers == 1:
%print "It is certain"
%
%elif answers == 2:
%print "Outlook good"
%
%elif answers == 3:
%print "You may rely on it"
%
%elif answers == 4:
%print "Ask again later"
%
%elif answers == 5:
%print "Concentrate and ask again"
%
%elif answers == 6:
%print "Reply hazy, try again"
%
%elif answers == 7:
%print "My reply is no"
%
%elif answers == 8:
%print "My sources say no"
%\end{lstlisting}
%\end{latin}
